% Section 4: the lowest common ancestor on the binary tree.
\section{Solving the tree: the lowest common ancestor}
\label{sec:lca}

\subsection{How the tree locates the route}

\begin{proposition}[Separation]
\label{prop:separation}
Let $R$ be a division with children $R_\ell,R_r$ and door $(u,v)$, $u\in R_\ell$,
$v\in R_r$. For cells $a,b$ in $R$, the route from $a$ to $b$ crosses $(u,v)$
if and only if $a$ and $b$ lie in different children.
\end{proposition}

\begin{proof}
$(u,v)$ is the only edge between the two halves (Proposition~\ref{prop:tree}).
If $a,b$ are on different sides, any route between them uses it. If they are on
the same side, that half's own spanning tree joins them, and the maze's unique
route is that one; it never leaves the half.
\end{proof}

% The separation lemma in one picture: the door of a division is on the route
% exactly when that division puts the two endpoints on opposite sides.
\begin{figure}[htbp]
  \centering
  \begin{subfigure}[b]{0.47\textwidth}\centering
    \begin{tikzpicture}[scale=0.72]
      \fill[cellgrey] (0,0) rectangle (8,5);
      \draw[black!55] (0,0) rectangle (8,5);
      \draw[wall segment] (4,0) -- (4,2);
      \draw[wall segment] (4,3) -- (4,5);
      \draw[door] (4,2) -- (4,3);
      \node[node vertex, fill=exitorange!25, draw=exitorange] (a) at (1.6,3.8) {};
      \node[node vertex, fill=exitorange!25, draw=exitorange] (b) at (6.4,1.2) {};
      \node[font=\tiny, above=1pt] at (a.north) {$a$};
      \node[font=\tiny, above=1pt] at (b.north) {$b$};
      \draw[linkblue, line width=1pt, rounded corners=3pt]
        (a) -- (2.6,2.5) -- (4,2.5) -- (5.4,2.5) -- (b);
      \node[font=\tiny, doorgreen, right] at (4.15,3.4) {the door};
    \end{tikzpicture}
    \caption{Opposite sides: the route \emph{must} cross, and the door is the
    only crossing there is. Recurse on both halves.}
    \label{fig:separated}
  \end{subfigure}\hfill
  \begin{subfigure}[b]{0.47\textwidth}\centering
    \begin{tikzpicture}[scale=0.72]
      \fill[cellgrey] (0,0) rectangle (8,5);
      \draw[black!55] (0,0) rectangle (8,5);
      \draw[wall segment] (4,0) -- (4,2);
      \draw[wall segment] (4,3) -- (4,5);
      \draw[door] (4,2) -- (4,3);
      \node[node vertex, fill=exitorange!25, draw=exitorange] (a) at (1.2,3.9) {};
      \node[node vertex, fill=exitorange!25, draw=exitorange] (b) at (2.9,1.1) {};
      \node[font=\tiny, above=1pt] at (a.north) {$a$};
      \node[font=\tiny, above=1pt] at (b.north) {$b$};
      \draw[linkblue, line width=1pt, rounded corners=3pt]
        (a) -- (2.2,3.0) -- (1.7,2.0) -- (b);
      \node[font=\tiny, black!55] at (6,2.5) {never entered};
    \end{tikzpicture}
    \caption{Same side: the route stays inside that half, and the other half is
    never looked at. Recurse once.}
    \label{fig:together}
  \end{subfigure}
  \caption{The separation lemma. Whether a division contributes a door to the
  route is decided by a containment test on its two children --- no cell of the
  maze is examined. The right-hand case is where the saving comes from: half the
  region is discarded in one step.}
  \label{fig:separation}
\end{figure}


Deciding which side a cell is on is a comparison of its coordinates against a
rectangle. So the search descends from the root while both cells stay on one
side: the first node that separates them is their \keyterm{lowest common
ancestor}, and its door is on the route. The route is then $a\to u$, the door,
$v\to b$ --- two smaller instances of the same problem, one per half.

\begin{algorithm}[H]
\caption{The route from $a$ to $b$ inside the tree node $R$.}
\label{alg:lca}
\SetKwFunction{Route}{Route}
\SetKwFunction{In}{In}
\Fn{\Route{$R$, $a$, $b$}}{
  \lIf{$R$ \textup{is a corridor}}{\Return the straight run $a\to b$}
  $(u,v)\leftarrow$ door of $R$\;
  \lIf{\In{$R_\ell$, $a$} \textup{\textbf{and}} \In{$R_\ell$, $b$}}{\Return \Route{$R_\ell$, $a$, $b$}}
  \lIf{\In{$R_r$, $a$} \textup{\textbf{and}} \In{$R_r$, $b$}}{\Return \Route{$R_r$, $a$, $b$}}
  \lIf{\In{$R_\ell$, $a$}}{\Return \Route{$R_\ell$, $a$, $u$} $\cdot$ \Route{$R_r$, $v$, $b$}}
  \Return \Route{$R_r$, $a$, $v$} $\cdot$ \Route{$R_\ell$, $u$, $b$}\;
}
\end{algorithm}

No cell of the maze is read. A corridor leaf returns its whole run in one step.

\subsection{Time complexity}

Two quantities, and they must be kept apart.

\keyterm{Locating the first door} is one descent to the lowest common ancestor:
$\Theta(\text{depth})$. The tree is $\Theta(\log n)$ deep in expectation ---
measured $2\log_2 n$ --- so this is $\Theta(\log n)$.

\keyterm{Producing the whole route} recurses on both sides of every door it
finds. At a division, either both cells stay on one side (one call) or they are
separated (two calls). With $p$ the probability of separation,
\[
  T(m)=(1+p)\,T(m/2)+O(1)
  \quad\Longrightarrow\quad
  T(n)=\Theta\!\left(n^{\log_2(1+p)}\right).
\]
Measured $p\approx0.61$ gives $\Theta(n^{0.68})$. It is not $\Theta(\log n)$
because $\log n$ is the cost of one descent and the search pays one per
separation; and it cannot be below $\Omega(d)$, the length of the route, since
the route itself is the output.

\begin{keypoint}
$\Theta(\log n)$ to find where the route crosses the first wall;
$\Theta(n^{0.68})$ to lay the whole route down. Both are far below A*'s
$\Theta(n)$ settled cells.
\end{keypoint}
